\documentclass{article}

\begin{document}

\textbf{Introduction}

\hfill

\textbf{BACKGROUND}

\hfill

Sir Isaac Newton brought to the world the idea of modeling the motion of
physical systems with equations. It was necessary to invent calculus along the
way, since fundamental equations of motion involve velocities and accelerations,
which are derivatives of position. His greatest single success was his discovery that
the motion of the planets and moons of the solar system resulted from a single
fundamental source: the gravitational attraction of the bodies. He demonstrated
that the observed motion of the planets could be explained by assuming that there
is a gravitational attraction between any two objects, a force that is proportional
to the product of masses and inversely proportional to the square of the distance
between them. The circular, elliptical, and parabolic orbits of astronomy were
no longer fundamental determinants of motion, but were approximations of laws
specified with differential equations. His methods are now used in modeling
motion and change in all areas of science.

Subsequent generations of scientists extended the method of using differential
equations to describe how physical systems evolve. But the method had
a limitation. While the differential equations were sufficient to determine the
behavior—in the sense that solutions of the equations did exist—it was frequently
difficult to figure out what that behavior would be. It was often impossible to write
down solutions in relatively simple algebraic expressions using a finite number of
terms. Series solutions involving infinite sums often would not converge beyond
some finite time.

When solutions could be found, they described very regular motion. Generations
of young scientists learned the sciences from textbooks filled with examples
of differential equations with regular solutions. If the solutions remained in
a bounded region of space, they settled down to either (A) a steady state, often
due to energy loss by friction, or (B) an oscillation that was either periodic or
quasiperiodic, akin to the clocklike motion of the moon and planets. (In the solar
system, there were obviously many different periods. The moon traveled around
the earth in a month, the earth around the sun in about a year, and Jupiter around
the sun in about 11.867 years. Such systems with multiple incommensurable
periods came to be called quasiperiodic.)

Scientists knew of systems which had more complicated behavior, such as
a pot of boiling water, or the molecules of air colliding in a room. However, since
these systems were composed of an immense number of interacting particles, the
complexity of their motions was not held to be surprising.

Around 1975, after three centuries of study, scientists in large numbers
around the world suddenly became aware that there is a third kind of motion, a
type (C) motion, that we now call “chaos”. The new motion is erratic, but not
simply quasiperiodic with a large number of periods, and not necessarily due to
a large number of interacting particles. It is a type of behavior that is possible in
very simple systems.

A small number of mathematicians and physicists were familiar with the
existence of a third type of motion prior to this time. James Clerk Maxwell, who
studied the motion of gas molecules in about 1860, was probably aware that even
a system composed of two colliding gas particles in a box would have neither
motion type A nor B, and that the long term behavior of the motions would for
all practical purposes be unpredictable. He was aware that very small changes
in the initial motion of the particles would result in immense changes in the
trajectories of the molecules, even if they were thought of as hard spheres.

Maxwell began his famous study of gas laws by investigating individual
collisions. Consider two atoms of equal mass, modeled as hard spheres. Give the
atoms equal but opposite velocities, and assume that their positions are selected
at random in a large three-dimensional region of space. Maxwell showed that if
they collide, all directions of travel will be equally likely after the collision. He
recognized that small changes in initial positions can result in large changes in
outcomes. In a discussion of free will, he suggested that it would be impossible
to test whether a leopard has free will, because one could never compute from a
study of its atoms what the leopard would do. But the chaos of its atoms is limited,
for, as he observed, “No leopard can change its spots!”

Henri Poincar´e in 1890 studied highly simplified solar systems of three
bodies and concluded that the motions were sometimes incredibly complicated.
(See Chapter 2). His techniques were applicable to a wide variety of physical
systems. Important further contributions were made by Birkhoff, Cartwright and
Littlewood, Levinson, Kolmogorov and his students, among others. By the 1960s,
there were groups of mathematicians, particularly in Berkeley and in Moscow,
striving to understand this third kind of motion that we now call chaos. But
only with the advent of personal computers, with screens capable of displaying
graphics, have scientists and engineers been able to see that important equations
in their own specialties had such solutions, at least for some ranges of parameters
that appear in the equations.

In the present day, scientists realize that chaotic behavior can be observed
in experiments and in computer models of behavior from all fields of science. The
key requirement is that the system involve a nonlinearity. It is now common for
experiments whose previous anomalous behavior was attributed to experiment
error or noise to be reevaluated for an explanation in these new terms. Taken
together, these new terms form a set of unifying principles, often called dynamical
systems theory, that cross many disciplinary boundaries.

The theory of dynamical systems describes phenomena that are common
to physical and biological systems throughout science. It has benefited greatly
from the collision of ideas from mathematics and these sciences. The goal of
scientists and applied mathematicians is to find nature’s unifying ideas or laws
and to fashion a language to describe these ideas. It is critical to the advancement
of science that exacting standards are applied to what is meant by knowledge.
Beautiful theories can be appreciated for their own sake, but science is a severe
taskmaster. Intriguing ideas are often rejected or ignored because they do not
meet the standards of what is knowledge.

The standards of mathematicians and scientists are rather different. Mathematicians
prove theorems. Scientists look at realisticmodels. Their approaches are
somewhat incompatible. The first papers showing chaotic behavior in computer
studies of very simple models were distasteful to both groups. The mathematicians
feared that nothing was proved so nothing was learned. Scientists said that models
without physical quantities like charge, mass, energy, or acceleration could not be
relevant to physical studies. But further reflection led to a change in viewpoints.
Mathematicians found that these computer studies could lead to new ideas that
slowly yielded new theorems. Scientists found that computer studies of much more
complicated models yielded behaviors similar to those of the simplistic models,
and that perhaps the simpler models captured the key phenomena.

Finally, laboratory experiments began to be carried out that showed unequivocal
evidence of unusual nonlinear effects and chaotic behavior in very
familiar settings. The new dynamical systems concepts showed up in macroscopic
systems such as fluids, common electronic circuits and low-energy lasers that were
previously thought to be fairly well understood using the classical paradigms. In
this sense, the chaotic revolution is quite different than that of relativity, which
shows its effects at high energies and velocities, and quantum theory, whose effects
are submicroscopic. Many demonstrations of chaotic behavior in experiments are
not far from the reader’s experience.

In this book we study this field that is the uncomfortable interface between
mathematics and science. We will look at many pictures produced by computers
and we try to make mathematical sense of them. For example, a computer study of
the driven pendulum in Chapter 2 reveals irregular, persistent, complex behavior
for ten million oscillations. Does this behavior persist for one billion oscillations?
The only way we can find out is to continue the computer study longer. However,
even if it continues its complex behavior throughout our computer study, we
cannot guarantee it would persist forever. Perhaps it stops abruptly after one
trillion oscillations; we do not know for certain. We can prove that there exist
initial positions and velocities of the pendulum that yield complex behavior
forever, but these choices are conceivably quite atypical. There are even simpler
models where we know that such chaotic behavior does persist forever. In this
world, pictures with uncertain messages remain the medium of inspiration.

There is a philosophy of modeling in which we study idealized systems
that have properties that can be closely approximated by physical systems. The
experimentalist takes the view that only quantities that can be measured have
meaning. Yet we can prove that there are beautiful structures that are so infinitely
intricate that they can never be seen experimentally. For example, we will see
immediately in Chapters 1 and 2 the way chaos develops as a physical parameter
like friction is varied. We see infinitely many periodic attractors appearing with
infinitely many periods. This topic is revisited in Chapter 12, where we show
how this rich bifurcation structure, called a cascade, exists with mathematical
certainty in many systems. This is a mathematical reality that underlies what
the experimentalist can see. We know that as the scientist finds ways to make
the study of a physical system increasingly tractable, more of this mathematical
structure will be revealed. It is there, but often hidden from view by the noise of
the universe. All science is of course dependent on simplistic models. If we study
a vibrating beam, we will generally not model the atoms of which it is made.
If we model the atoms, we will probably not reflect in our model the fact that
the universe has a finite age and that the beam did not exist for all time. And
we do not include in our model (usually) the tidal effects of the stars and the
planets on our vibrating beam. We ignore all these effects so that we can isolate
the implications of a very limited list of concepts.

It is our goal to give an introduction to some of the most intriguing ideas in
dynamics, the ideas we love most. Just as chemistry has its elements and physics
has its elementary particles, dynamics has its fundamental elements: with names
like attractors, basins, saddles, homoclinic points, cascades, and horseshoes. The
ideas in this field are not transparent. As a reader, your ability to work with these
ideas will come from your own effort.We will consider our job to be accomplished
if we can help you learn what to look for in your own studies of dynamical systems
of the world and universe.

\hfill

\textbf{ABOUT T H E BOOK}

\hfill

As we developed the drafts of this book, we taught six one semester classes at
George Mason University and the University of Maryland. The level is aimed at
undergraduates and beginning graduate students. Typically, we have used parts
of Chapters 1–9 as the core of such a course, spending roughly equal amounts of
time on iterated maps (Chapters 1–6) and differential equations (Chapters 7–9).
Some of the maps we use as examples in the early chapters come from differential
equations, so that their importance in the subject is stressed. The topics of stable
manifolds, bifurcations, and cascades are introduced in the first two chapters and
then developed more fully in the Chapters 10, 11, and 12, respectively. Chapter
13 on time series may be profitably read immediately after Chapter 4 on fractals,
although the concepts of periodic orbit (of a differential equation) and chaotic
attractor will not yet have been formally defined.

The impetus for advances in dynamical systems has come from many
sources: mathematics, theoretical science, computer simulation, and experimenix
tal science. We have tried to put this book together in a way that would reflect
its wide range of influences.

We present elaborate dissections of the proofs of three deep and important
theorems: The Poincar´e-Bendixson Theorem, the Stable Manifold Theorem, and
the Cascade Theorem. Our hope is that including them in this form tempts you
to work through the nitty-gritty details, toward mastery of the building blocks as
well as an appreciation of the completed edifice.
Additionally, each chapter contains a special feature called a Challenge,
in which other famous ideas from dynamics have been divided into a number
of steps with helpful hints. The Challenges tackle subjects from period-three
implies chaos, the cat map, and Sharkovskii’s ordering through synchronization
and renormalization.We apologize in advance for the hints we have given, when
they are of no help or even mislead you; for one person’s hint can be another’s
distraction.

The Computer Experiments are designed to present you with opportunities
to explore dynamics through computer simulation, the venue through which
many of these concepts were first discovered. In each, you are asked to design
and carry out a calculation relevant to an aspect of the dynamics. Virtually all
can be successfully approached with a minimal knowledge of some scientific
programming language. Appendix B provides an introduction to the solution of
differential equations by approximate means, which is necessary for some of the
later Computer Experiments.

If you prefer not to work the Computer Experiments from scratch, your
task can be greatly simplified by using existing software. Several packages
are available. Dynamics: Numerical Explorations by H.E. Nusse and J.A. Yorke
(Springer-Verlag 1994) is the result of programs developed at the University of
Maryland. Dynamics, which includes software for Unix and PC environments,
was used to make many of the pictures in this book. The web site for Dynamics
is www.ipst.umd.edu/dynamics. We can also recommend Differential and
Difference Equations through Computer Experiments by H. Kocak (Springer-Verlag,
1989) for personal computers. A sophisticated package designed for Unix platforms
is dstool, developed by J. Guckenheimer and his group at Cornell University.
In the absence of special purpose software, general purpose scientific computing
environments such as Matlab, Maple, and Mathematica will do nicely.

The Lab Visits are short reports on carefully selected laboratory experiments
that show how the mathematical concepts of dynamical systems manifest
themselves in real phenomena.We try to impart some flavor of the setting of the
experiment and the considerable expertise and care necessary to tease a new secret
from nature. In virtually every case, the experimenters’ findings far surpassed
what we survey in the Lab Visit.We urge you to pursue more accurate and detailed
discussions of these experiments by going straight to the original sources.



\end{document}


